\section{Practical interpretation of scaled errors and intervals}\label{sec:intervals}
The primary candidate has MASE 1.004--1.110, despite relative MAE below one against contemporaneous persistence (Table~\ref{tab:release_scaled}). These measures ask different questions: MASE uses the entity's training-period naive scale, whereas relative MAE compares two evaluation-period forecasts on the same rows. The denominators, aggregation and even retained row populations differ. A lower pooled rupee MAE therefore supports a modest historical average gain, not a universal victory over all benchmarks, every market or the training-period scale. MASE exclusions (4.4--4.5\%) and subgroup failures remain part of the result, not inconvenient observations to omit.

\begin{table}[t]\centering\small\setlength{\tabcolsep}{4pt}
\caption{Primary study scaled errors and observed nominal 80\% price-interval performance. C/P denote candidate/persistence. Width is divided by origin price. MASE excludes unusable training scales.}\label{tab:release_scaled}
\begin{tabular}{@{}llrrrr@{}}\toprule
Regime & Horizon & MASE C & MASE P & Coverage 80\% & Rel. width\\\midrule
cleaned & next & 1.100 & 1.156 & 77.7\% & 0.217 \\
cleaned & h3 & 1.004 & 1.060 & 79.1\% & 0.272 \\
cleaned & h7 & 1.074 & 1.109 & 79.4\% & 0.338 \\
cleaned & h14 & 1.044 & 1.061 & 79.4\% & 0.442 \\
unfiltered & next & 1.110 & 1.168 & 77.8\% & 0.219 \\
unfiltered & h3 & 1.014 & 1.071 & 79.1\% & 0.273 \\
unfiltered & h7 & 1.080 & 1.117 & 79.5\% & 0.342 \\
unfiltered & h14 & 1.048 & 1.065 & 79.5\% & 0.443 \\
\bottomrule\end{tabular}

\end{table}

Observed price-interval coverage is 77.7--79.5\%, compared with nominal 80\%; undercoverage is most evident at the next report. This measured historical frequency does not validate future coverage under dependence or changing conditions. Wider intervals do not guarantee an economically useful decision. These price intervals must also be distinguished from block-bootstrap sensitivity intervals for the aggregate MAE ratio: the latter describe conditional resampling variability, not ranges for a future mandi price. Neither is a profit guarantee.
